\chapter{How to Read a Chinese Poem}

\section{Course overview}

This course reads classical Chinese poetry (\zht{詩} \emph{shī}, and later the lyric
\zht{詞} \emph{cí}) from the \emph{Book of Songs} (c.\ 1000--600 BCE) to the Song dynasty
(960--1279), with a brief look at what came after. Everything is read in English
translation alongside the originals: you do not need to read Chinese, but you need to
be able to \emph{look} at Chinese.
By the end of term we should be able to count characters in a line, spot a rhyme, find a
parallel couplet and say what a translator has gained or lost.

\begin{important}[Assessment (from the syllabus)]
\begin{itemize}
  \item Weekly reading responses (short, one poem each): 20\%.
  \item Two ``close reading'' papers: take one short poem, compare two translations, and
        argue about what the original does that the translations cannot.
  \item Midterm (identifications: poet, dynasty, form, one key image) and final exam
        (identifications + essay).
  \item Translation exercise: produce your own version of a short poem from a
        character-by-character gloss and defend your choices.
\end{itemize}
\end{important}

\section{A timeline to hang everything on}

Chinese literary history is conventionally organised by dynasty. We will not need every
dynasty, but we need these.

\begin{center}
\begin{tikzpicture}[x=0.0062cm, y=1cm, every node/.style={font=\scriptsize}]
  \draw[thick, -{Stealth}] (-1100,0) -- (1400,0);
  \foreach \yr in {-1000,-500,0,500,1000} {
    \draw (\yr,0.1) -- (\yr,-0.1) node[below] {\yr};
  }
  % dynasty bars
  \fill[orange!30] (-1046,0.25) rectangle (-256,0.65);
  \node at (-650,0.45) {Zhou \zht{周}};
  \fill[red!25] (-206,0.25) rectangle (220,0.65);
  \node at (7,0.45) {Han \zht{漢}};
  \fill[gray!25] (220,0.25) rectangle (589,0.65);
  \node at (405,0.45) {Six Dyn.};
  \fill[blue!25] (618,0.25) rectangle (907,0.65);
  \node at (762,0.45) {Tang \zht{唐}};
  \fill[green!25] (960,0.25) rectangle (1279,0.65);
  \node at (1120,0.45) {Song \zht{宋}};
  % works and poets
  \draw[gray] (-800,0.65) -- (-800,1.3) node[above] {\emph{Shijing} \zht{詩經}};
  \draw[gray] (-300,0.65) -- (-300,1.9) node[above] {Qu Yuan \zht{屈原}};
  \draw[gray] (150,0.65) -- (150,1.3) node[above] {\emph{Nineteen Old Poems}};
  \draw[gray] (400,0.65) -- (400,1.9) node[above] {Tao Yuanming \zht{陶淵明}};
  \draw[gray] (730,0.65) -- (730,3.0) node[above] {Wang Wei, Li Bai, Du Fu};
  \draw[gray] (840,0.65) -- (840,2.5) node[above right, xshift=-6pt] {Bai Juyi, Li Shangyin};
  \draw[gray] (1090,0.65) -- (1090,1.9) node[above] {Su Shi, Li Qingzhao};
  \node[text=red!60!black] at (755,-0.75) {An Lushan rebellion 755};
  \draw[red!60!black, thick] (755,-0.1) -- (755,-0.55);
\end{tikzpicture}
\end{center}

\begin{remark}
``Six Dynasties'' (\zht{六朝}, 220--589) is a convenient name for the period of
division between the Han and the reunification under the Sui (581--618). Tao Yuanming
lives here. The Tang is the ``golden age'' of \emph{shī}; the Song is the golden age of
\emph{cí}.
\end{remark}

\section{Four features of the language}

\subsection{Monosyllables and the character}

Each Chinese character (\zht{字} \emph{zì}) is written as one block and read as one
syllable, and in Classical Chinese most characters are whole words. So ``a five-character
line'' (\zht{五言} \emph{wǔyán}) is literally five syllables and usually five words. A
twenty-character quatrain can hold an entire scene. English needs articles, prepositions,
tense markers and pronouns that the original simply does not have.

\begin{example}[Counting]
\zht{床前明月光} \emph{chuáng qián míng yuè guāng} = ``bed / front / bright / moon /
light''. Five characters, five syllables, no verb, no ``I'', no ``the''. Any English
version has to decide whose bed, what time, and whether the light is \emph{on} the floor
or merely \emph{in front of} the bed.
\end{example}

\subsection{What the grammar leaves open}

Classical poetic language routinely omits:
\begin{itemize}
  \item \textbf{the subject}: who is looking at the moon is never said;
  \item \textbf{tense}: there is no grammatical past/present, only context;
  \item \textbf{number}: \zht{山} \emph{shān} can be ``a mountain'' or ``mountains'';
  \item \textbf{connectives}: lines are juxtaposed without ``because'', ``while'', ``so''.
\end{itemize}
This openness is not vagueness for its own sake. It lets the poem present a scene
\emph{as it is experienced}, before the experiencer is separated from it. Translators into
English must constantly decide what the poem chose not to decide.

\begin{aside}
The ``Imagist'' poets of the 1910s (Ezra Pound's \emph{Cathay}, 1915) were excited
precisely by this juxtaposition of images without connectives. Pound worked from Ernest
Fenollosa's notes and did not read Chinese; his versions are wonderful English poems and
often unreliable translations. We will come back to this tension all term.
\end{aside}

\subsection{Imagery: the image carries the feeling}

Classical poets rarely announce emotions outright; they set out images whose conventional
associations do the work. A short list of images we will meet constantly:

\begin{center}
\begin{tabular}{lll}
\toprule
Image & Pinyin & Conventional association \\
\midrule
\zht{月} moon & yuè & distance, separated people looking at the same moon \\
\zht{柳} willow & liǔ & parting (\zht{柳} sounds like \zht{留} ``stay''; branches given at farewells) \\
\zht{雁} wild goose & yàn & letters, news from far away, autumn \\
\zht{秋} autumn & qiū & decline, ageing, the scholar's melancholy \\
\zht{酒} wine & jiǔ & conviviality; also escape, carpe diem \\
\zht{浮雲} drifting clouds & fúyún & the wanderer; things that obscure (slanderers) \\
\zht{菊} chrysanthemum & jú & reclusion, Tao Yuanming, endurance into autumn \\
\zht{猿} gibbon & yuán & its cry in the Yangzi gorges makes travellers weep \\
\bottomrule
\end{tabular}
\end{center}

\subsection{Allusion}

An \emph{allusion} (\zht{典故} \emph{diǎngù}) points to an earlier text or story. A
learned reader hears two poems at once. Allusion becomes denser over time: Tao Yuanming
alludes to the classics, Du Fu to everything, and Li Shangyin's allusions are sometimes the
whole puzzle.

\section{Parallelism}

\begin{definition}[Parallel couplet]
A \emph{parallel couplet} (\zht{對仗} \emph{duìzhàng}) is a pair of lines in which each
position in the first line matches the same position in the second by grammatical
category (noun with noun, verb with verb, number with number, colour with colour) and
usually by meaning (contrast or complement).
\end{definition}

The classic small example is the first couplet of Wang Zhihuan's (\zht{王之渙}, 688--742)
``Climbing Stork Tower'' (\zht{登鸛雀樓}):

\begin{center}
\begin{tikzpicture}[ch/.style={draw, minimum size=0.85cm, font=\large},
  gl/.style={font=\scriptsize, text width=1.3cm, align=center}]
  \foreach \c [count=\i] in {白,日,依,山,盡} {
    \node[ch] (a\i) at (1.5*\i,1.6) {\zht{\c}};
  }
  \foreach \c [count=\i] in {黃,河,入,海,流} {
    \node[ch] (b\i) at (1.5*\i,0) {\zht{\c}};
  }
  \foreach \i in {1,...,5} { \draw[blue!60, thick, <->] (a\i.south) -- (b\i.north); }
  \node[gl] at (1.5,2.35) {white};   \node[gl] at (3,2.35) {sun};
  \node[gl] at (4.5,2.35) {lean on}; \node[gl] at (6,2.35) {mountain};
  \node[gl] at (7.5,2.35) {end};
  \node[gl] at (1.5,-0.75) {yellow}; \node[gl] at (3,-0.75) {river};
  \node[gl] at (4.5,-0.75) {enter};  \node[gl] at (6,-0.75) {sea};
  \node[gl] at (7.5,-0.75) {flow};
  \node[font=\scriptsize, text=blue!60!black, anchor=west] at (8.3,0.8)
    {colour $\cdot$ noun $\cdot$ verb $\cdot$ noun $\cdot$ verb};
\end{tikzpicture}
\end{center}

\begin{center}
\begin{tabular}{p{0.32\textwidth}p{0.58\textwidth}}
\toprule
\zht{白日依山盡，} & The white sun, leaning on the hills, is gone; \\
\zht{黃河入海流。} & the Yellow River flows on into the sea. \\
\zht{欲窮千里目，} & If you would see to the end of a thousand miles, \\
\zht{更上一層樓。} & climb one more storey of the tower. \\
\bottomrule
\end{tabular}
\end{center}

Note how the parallel couplet sets up a cosmos: west (mountains, sunset) against east
(the river running to the sea), sky against earth, an ending against a continuing. The
second couplet is \emph{also} parallel (``wish / exhaust / thousand / \emph{li} / eye'' and
``further / ascend / one / storey / tower'': number against number), but it reads as a
single flowing sentence. This ``flowing-water pair'' (\zht{流水對}) was much admired.

\section{The problem of translation}

\subsection{Choices every translator faces}
\begin{enumerate}
  \item \textbf{Line length.} Five-character lines become English lines of 8--12 words.
        Do you keep them short (and choppy) or let them breathe?
  \item \textbf{Rhyme.} Nearly every Chinese poem rhymes. English rhyme often forces
        padding; most modern translators drop it.
  \item \textbf{Pronouns and tense.} Must be added. ``I'' makes the poem personal; leaving
        the speaker out makes it impersonal and imagistic.
  \item \textbf{Allusions.} Explain in the line, in a footnote, or not at all?
  \item \textbf{Proper names.} ``Yellow River'' (translated) or ``the Huang He''
        (transliterated)?
\end{enumerate}

\subsection{Case study: Li Bai's ``Quiet Night Thoughts''}

\begin{center}
\begin{tabular}{p{0.25\textwidth}p{0.30\textwidth}p{0.35\textwidth}}
\toprule
Original & Pinyin & Character gloss \\
\midrule
\zht{床前明月光，} & chuáng qián míng yuè guāng & bed front bright moon light \\
\zht{疑是地上霜。} & yí shì dì shàng shuāng & suspect is ground on frost \\
\zht{舉頭望明月，} & jǔ tóu wàng míng yuè & raise head gaze bright moon \\
\zht{低頭思故鄉。} & dī tóu sī gù xiāng & lower head think old village \\
\bottomrule
\end{tabular}
\end{center}

Here are two versions I wrote for the first close-reading exercise, deliberately at
opposite ends.

\begin{example}[Version A: literal, keeping the line count]
Before my bed, the bright moon's light: \\
I took it for frost upon the ground. \\
I lift my head and gaze at the bright moon; \\
I lower my head and think of home.
\end{example}

\begin{example}[Version B: free, imagistic]
Moonlight by the bed ---\\
frost on the floor?\\
Head up: the moon.\\
Head down: home.
\end{example}

\begin{important}[What each version gains and loses]
\begin{itemize}
  \item A keeps the four-line shape and the repeated \zht{明月} ``bright moon'' (lines 1
        and 3), and the strict parallel of lines 3--4 (\zht{舉頭} / \zht{低頭}, ``raise
        head'' / ``lower head''). But it adds ``my'' and ``I'' four times, which the
        original never says, and the past tense ``took'' is a guess.
  \item B keeps the speed and impersonality and makes the head-up / head-down hinge very
        visible. But it loses \zht{疑} ``to suspect, to be in doubt'' (turned into a
        question mark), loses \zht{思} ``to think of, to long for'', and its brevity
        sounds like a modern haiku rather than a Tang song.
  \item Neither can reproduce the rhyme \emph{guāng / shuāng / xiāng}, nor the fact that
        the poem is made of the plainest possible words: every Chinese schoolchild knows
        it.
\end{itemize}
\end{important}

\begin{aside}
The text above is the version taught in schools since the Ming dynasty. Song-dynasty
editions read \zht{床前看月光} ``before my bed I look at the moonlight'' and \zht{舉頭望山月}
``raise my head to gaze at the mountain moon''. Also, \zht{床} may not be a bed at all:
some scholars argue it is a folding seat, or the railing round a well. Even the most famous
poem in the language has textual and lexical uncertainty.
\end{aside}

\section{A procedure for close reading}

When you meet a new poem, work through these steps (this became the structure of every
reading response):
\begin{enumerate}
  \item \textbf{Count}: characters per line, lines per poem. Five or seven? Four or eight?
        Irregular?
  \item \textbf{Listen}: which line-final characters rhyme? (Read the pinyin aloud, but
        remember that modern Mandarin is not Tang pronunciation.)
  \item \textbf{Map}: which couplets are parallel? Draw the lines between matched words.
  \item \textbf{Picture}: list the images. What are their conventional associations?
  \item \textbf{Locate}: when, where, under what circumstances? Is there a title that
        tells you the occasion (``Seeing off \ldots'', ``Climbing \ldots'')?
  \item \textbf{Turn}: where does the poem shift, usually in the last couplet, from scene
        (\zht{景} \emph{jǐng}) to feeling (\zht{情} \emph{qíng})?
\end{enumerate}

\section*{Summary}
\begin{itemize}
  \item One character = one syllable = (usually) one word; lines are counted in
        characters (five- and seven-character lines dominate).
  \item The poetic language omits subject, tense, number and connectives; translators must
        supply them, and every choice is an interpretation.
  \item Images carry conventional emotional associations (moon, willow, goose, autumn).
  \item Parallel couplets match position by position in grammar and meaning.
  \item Close reading: count, listen, map, picture, locate, find the turn from scene to
        feeling.
\end{itemize}
